\subsection{Inference and Learning in Graphical Models}

Inference in a graphical model computes marginal, conditional, or most likely latent quantities from the joint law encoded by the graph.

\begin{definition} \label{def:sl-graphical-model-inference-tasks} \hlt{Inference Tasks}\\
Common tasks include computing marginals $p(x_A)$, conditionals $p(x_A\mid x_B)$, maximum a posteriori assignments, and posterior distributions over latent variables.
\end{definition}

\begin{method} \label{method:sl-message-passing} \hlt{Message Passing}\\
On trees, belief propagation sends local messages along edges and computes exact marginals.  On loopy graphs, analogous iterations often serve as approximations.
\end{method}

\begin{definition} \label{def:sl-variational-inference} \hlt{Variational Inference}\\
Variational inference approximates an intractable posterior $p(z\mid x)$ by a tractable distribution $q(z)$, commonly by minimizing $\operatorname{KL}(q(z)\Vert p(z\mid x))$ or maximizing an evidence lower bound.
\end{definition}

\begin{remark} \label{rem:sl-graphical-learning} \hlt{Learning Parameters and Structure}\\
Parameter learning uses likelihood or regularized likelihood when the graph is fixed.  Structure learning searches over graphs and therefore combines statistical scoring with combinatorial optimization.
\end{remark}
